\section*{Capítulo \ref{ch:b1:crystals-ionic-covalent} --- Cristales II: sólidos iónicos, covalentes y moleculares}

\begin{solution}{exo:b1:crystals-ionic-covalent:1}
Aniones: $8 \times \frac18 + 6 \times \frac12 = 4$; cationes: $12 \times
\frac14 + 1 = 4$. Cada ion tiene 6 vecinos del otro tipo (6:6). Cuatro
cationes y cuatro aniones: \ce{NaCl}.
\end{solution}

\begin{solution}{exo:b1:crystals-ionic-covalent:2}
La mitad de la diagonal principal: $a\sqrt3/2 = 412.3 \times 0.866 = \qty{357.1}{pm}$.
Un \ce{Cs+} y $8 \times \frac18 = 1$ \ce{Cl-}: una unidad fórmula.
\end{solution}

\begin{solution}{exo:b1:crystals-ionic-covalent:3}
$x = 102/181 = 0.56$, entre 0.414 y 0.732: coordinación octaédrica (6),
el \omterm{def:b1:crystals-ionic-covalent:structure-type}{tipo sal gema}.
\end{solution}

\begin{solution}{exo:b1:crystals-ionic-covalent:4}
Metálico: el cobre. Iónicos: el cloruro de sodio y el fluoruro de calcio.
Covalentes: el diamante y el cuarzo. Moleculares: el diyodo y el hielo.
\end{solution}

\begin{solution}{exo:b1:crystals-ionic-covalent:5}
A lo largo de una arista, anión, catión y anión en contacto:
$a = 2(r_+ + r_-)$. Sobre una diagonal de cara hay dos aniones a la
distancia $a\sqrt2/2$, que no deben solaparse: $a\sqrt2/2 \ge 2r_-$, de
modo que $(r_+ + r_-)\sqrt2 \ge 2r_-$, es decir,
$r_+/r_- \ge \sqrt2 - 1$.
\end{solution}

\begin{solution}{exo:b1:crystals-ionic-covalent:6}
$\rho = 4 \times 58.44/(6.022 \times 10^{23} \times (5.641 \times
10^{-8})^3) = \qty{2.16}{g/cm^3}$, frente a \qty{2.17}{g/cm^3} medida.
\end{solution}

\begin{solution}{exo:b1:crystals-ionic-covalent:7}
$\rho = 168.36/(6.022 \times 10^{23} \times (4.123 \times 10^{-8})^3) =
\qty{3.99}{g/cm^3}$, el valor medido.
\end{solution}

\begin{solution}{exo:b1:crystals-ionic-covalent:8}
\ce{Zn-S} $= a\sqrt3/4 = \qty{234.2}{pm}$;
$\rho = 4 \times 97.44/(6.022 \times 10^{23} \times (5.409 \times
10^{-8})^3) = \qty{4.09}{g/cm^3}$, próxima a los \qty{4.04}{g/cm^3}
medidos en la esfalerita natural.
\end{solution}

\begin{solution}{exo:b1:crystals-ionic-covalent:9}
\ce{C-C} $= 245.6/\sqrt3 = \qty{141.8}{pm}$; entre capas, $669.6/2 =
\qty{334.8}{pm}$. Volumen de la celda: $\frac{\sqrt3}{2}a^2c = \qty{3.498e-23}{cm^3}$,
$\rho = 4 \times 12.01/(6.022 \times 10^{23} \times 3.498 \times 10^{-23})
= \qty{2.28}{g/cm^3}$, frente a \qty{3.52}{g/cm^3} para el diamante: las
capas están muy separadas, unidas solo por fuerzas de London.
\end{solution}

\begin{solution}{exo:b1:crystals-ionic-covalent:10}
Véase la demostración de la \cref{prop:b1:crystals-ionic-covalent:diamond}:
$C = \pi\sqrt3/16 = 0.34$. La dureza procede de la fuerza y de la rigidez
tridimensional de los \omterm{def:b1:lewis-resonance-vsepr:lewis}{enlaces covalentes}, no del empaquetamiento: cada
átomo solo tiene cuatro vecinos, a ángulos fijos, lo que deja mucho
espacio vacío.
\end{solution}

\begin{solution}{exo:b1:crystals-ionic-covalent:11}
\ce{Si-Si} $= a\sqrt3/4 = \qty{235.2}{pm}$, \omterm{def:b1:periodicity:radius}{radio covalente}
\qty{117.6}{pm}; $\rho = 8 \times 28.09/(6.022 \times 10^{23} \times
(5.431 \times 10^{-8})^3) = \qty{2.33}{g/cm^3}$, el valor medido.
\end{solution}

\begin{solution}{exo:b1:crystals-ionic-covalent:12}
$x = 152/181 = 0.84 > 0.732$: la regla predice el \omterm{def:b1:crystals-ionic-covalent:structure-type}{tipo cloruro de cesio}.
En la estructura de la sal gema, \ce{Rb-Cl} $= a/2 = \qty{329.1}{pm}$,
próxima a $152 + 181 = \qty{333}{pm}$: la estructura observada es
coherente. La regla trata los iones como esferas rígidas y no tiene en
cuenta su \omterm{def:b1:periodicity:polarisability}{polarizabilidad} ni la pequeña diferencia de energía entre las
dos estructuras (el cloruro de rubidio adopta el \omterm{def:b1:crystals-ionic-covalent:structure-type}{tipo cloruro de cesio} a
alta presión); es una guía, no una ley.
\end{solution}

\begin{solution}{pb:b1:crystals-ionic-covalent:1}
\textbf{1.} En los 8 vértices y en los 6 centros de las caras: $8 \times \frac18 + 6
\times \frac12 = 4$ iones.
\textbf{2.} 8 \omterm{def:b1:crystals-metals:site}{huecos tetraédricos}, todos ocupados: 8 \ce{F-} para 4
\ce{Ca^{2+}}, proporción 1:2, \ce{CaF2}.
\textbf{3.} \ce{F-}: 4 \ce{Ca^{2+}} en los vértices de un tetraedro.
\ce{Ca^{2+}}: 8 \ce{F-} en los vértices de un cubo.
\textbf{4.} $a\sqrt3/4 = 546.3 \times 0.4330 = \qty{236.6}{pm}$.
\textbf{5.} $112 + 131 = \qty{243}{pm}$: dentro de un 3\,\%.
\textbf{6.} Los iones \ce{F-} forman una disposición cúbica simple de
arista $a/2 =
\qty{273.2}{pm}$, algo más que $2 \times 131 = \qty{262}{pm}$: casi se
tocan.
\textbf{7.} $x = 112/131 = 0.85$.
\textbf{8.} $x > 0.732$: coordinación 8.
\textbf{9.} El \omterm{def:b1:crystals-ionic-covalent:structure-type}{tipo cloruro de cesio} tiene tantos cationes como aniones;
con el doble de aniones, solo la mitad de los cubos de aniones pueden
alojar un catión.
\textbf{10.} Los 8 \ce{F-} de la celda forman un cubo de 8 cubos pequeños
de arista $a/2$; los iones \ce{Ca^{2+}} ocupan los centros de 4 de ellos,
alternados, como las casillas negras de un tablero de ajedrez
tridimensional.
\textbf{11.} $C = \frac43\pi(4 \times 112^3 + 8 \times 131^3)/546.3^3 = 0.61$.
\textbf{12.} \ce{O^{2-}} en la \omterm{def:b1:crystals-metals:lattice}{red} FCC (vértices y centros de las caras),
y \ce{Li+} en los 8 \omterm{def:b1:crystals-metals:site}{huecos tetraédricos}.
\textbf{13.} $a\sqrt3/4 = \qty{200.0}{pm}$; $59 + 142 = \qty{201}{pm}$.
\textbf{14.} $\rho = 4 \times 29.88/(6.022 \times 10^{23} \times (4.619
\times 10^{-8})^3) = \qty{2.01}{g/cm^3}$, el valor medido.
\textbf{15.} La blenda: el diamante es una blenda con carbono en los dos
tipos de posiciones.
\textbf{16.} $a\sqrt3/4 = \qty{154.5}{pm}$; $\rho = 8 \times
12.01/(6.022 \times 10^{23} \times (3.567 \times 10^{-8})^3) =
\qty{3.52}{g/cm^3}$.
\textbf{17.} En el diamante solo está ocupada la mitad de los \omterm{def:b1:crystals-metals:site}{huecos
tetraédricos}, y por átomos tan grandes como los de la \omterm{def:b1:crystals-metals:lattice}{red}, que solo se
tocan a lo largo de los enlaces: $C = 0.34$. En la fluorita todos los
huecos están ocupados y los cationes pequeños llenan los espacios entre
aniones grandes: $C = 0.61$.
\textbf{18.} $40.08 + 2 \times 19.00 = \qty{78.08}{g/mol}$.
\textbf{19.} $4 \times 78.08/(6.022 \times 10^{23}) = \qty{5.186e-22}{g}$.
\textbf{20.} $(5.463 \times 10^{-8})^3 = \qty{1.630e-22}{cm^3}$.
\textbf{21.} Un \ce{F-} que falta cada cien celdas quita $19.00/(100
\times 312.3)$ de la masa, un 0.06\,\%: invisible con tres cifras (y
también debe faltar un \ce{Ca^{2+}}, o el \omterm{def:b1:crystals-metals:crystal}{cristal} no sería neutro).
\textbf{22.} $\rho = 5.186 \times 10^{-22}/1.630 \times 10^{-22} =
\textbf{\qty{3.18}{g/cm^3}}$, exactamente la densidad medida: el modelo de
la celda, construido a partir de una sola medida de $a$ por rayos X,
explica la masa del mineral.
\end{solution}
